\section{Preliminary GNN Encoder Viability}
\label{subsec:gnn-viability}

Before screening individual graph-design choices, an exploratory study checks
that a graph actor can match the tuned MLP baseline. This is a prerequisite for
transfer: the agents must be able to share one size-independent encoder without
sacrificing survival.

\paragraph{Question.}
Can a shared graph encoder perform as well as the MLP? If each agent requires a
separate encoder, or if the GNN cannot train stably, its structural
transferability is lost.

\paragraph{Baseline and changes.}
The study compares several GINE variants. It varies encoder size (light or
heavy), actor sharing (shared or unshared encoder), critic architecture (MLP or
GNN), and the training optimizations established for the MLP baseline,
including optimized critic updates and reward normalization.

\paragraph{Protocol.}
All GINE runs use \texttt{gnn\_type = gine}, $15{,}000{,}000$ total
timesteps, evaluation every $80{,}000$ timesteps, $2{,}000$ rollout steps,
deterministic evaluation, and the same 80/20 chronic split used by the MLP
experiments. GINE is used because it incorporates edge features and is
expressive over graph structures~\cite{xu2018how}. The complete graph-actor
pipeline is described in Appendix~\ref{app:gine-architecture};
Appendix~\ref{app:experiment-configurations} gives the full configurations in
Tables~\ref{tab:gine-run-configurations} and~\ref{tab:heavy-gine-configurations}.

\paragraph{Evidence.}
Table~\ref{tab:gine-rerun-summary} and Figure~\ref{fig:rerun-gine-best00}
compare the main architectural improvements with the \texttt{MLP Baseline}.
Figure~\ref{fig:gine-best-search2} shows the broader search over isolated
changes.

\begin{table}[H]
  \centering
  \small
  \caption{GINE rerun survival summary.}
  \label{tab:gine-rerun-summary}
  \resizebox{\textwidth}{!}{%
  \begin{tabular}{llrr}
    \toprule
    Run & Intended change & Final survival (\%) & Peak survival (\%) \\
    \midrule
    \texttt{Heavy GINE Baseline} & Historical GINE baseline & 59.2 & 93.9 \\
    \texttt{Light GINE (Concat, MLP Critic)} & Light GINE with MLP critic & 78.9 & 90.4 \\
    \texttt{No Bias, Opt Critic} & Optimized critic and bias 0.0 & 82.5 & 85.7 \\
    \texttt{Optimized Light GINE} & Light optimized GINE with MLP critic & 95.7 & 97.4 \\
    \bottomrule
  \end{tabular}
  }
\end{table}

\begin{figure}[H]
  \centering
  \includegraphics[width=\textwidth]{figures/rerun_gine_best00_vs_best10_14.png}
  \caption{Seed-aggregated test episodic survival comparing the
  \texttt{MLP Baseline} with high-performing GINE architectural variants.}
  \label{fig:rerun-gine-best00}
\end{figure}

\begin{figure}[p]
  \centering
  \includegraphics[width=\textwidth,height=0.9\textheight,keepaspectratio]{figures/gine_best_search2_baseline_comparisons.png}
  \caption{GINE survival comparisons. Each subplot compares the
  \texttt{MLP Baseline} with one variant from
  \texttt{configs/gine\_s0\_s1\_s2}. Full hyperparameters and architectural
  details are given in Appendices~\ref{app:gine-architecture}
  and~\ref{app:experiment-configurations}.}
  \label{fig:gine-best-search2}
\end{figure}

\paragraph{Interpretation.}
The broad search gives three conclusions:
\begin{itemize}
  \item \textbf{Shared encoder viability.} The unshared-actor variant does not
  outperform the shared variants, so one encoder can serve all agents without
  becoming the performance bottleneck.
  \item \textbf{Critic architecture.} A GNN centralized critic is unstable and
  degrades performance. The strongest variants pair the shared GNN actor with
  the standard MLP critic.
  \item \textbf{Optimization dependence.} Optimized critic updates, reward
  normalization, and removal of the initial do-nothing bias are also important
  for stable GNN training.
\end{itemize}

\noindent The strongest result is the \texttt{Optimized Light GINE}, which
combines the shared light GINE actor, the MLP critic, and the baseline
optimizations to reach $95.7\%$ final survival. A shared GNN is therefore
competitive on \texttt{bus14}, justifying the representation screens that
follow.
